\documentclass[11pt]{article}
\usepackage{amsmath,amssymb,amsthm}
\usepackage{tikz}
\usepackage{hyperref}
\usepackage[margin=1in]{geometry}

\newcommand{\N}{\mathbb{N}}
\newcommand{\Z}{\mathbb{Z}}
\newcommand{\Q}{\mathbb{Q}}
\newcommand{\GK}{\operatorname{GK}}
\newcommand{\CPP}{\operatorname{CPP}}
\newcommand{\ALCD}{\operatorname{ALCD}}
\newcommand{\qpoch}[2]{(#1;#2)_\infty}
\newcommand{\qpochn}[3]{(#1;#2)_{#3}}
\newcommand{\qbin}[2]{\genfrac{[}{]}{0pt}{}{#1}{#2}}
\DeclareMathOperator{\tr}{tr}
\DeclareMathOperator{\wt}{wt}
\DeclareMathOperator{\sgn}{sgn}
\DeclareMathOperator{\Fix}{Fix}
\DeclareMathOperator{\id}{id}
\DeclareMathOperator{\supp}{supp}

\theoremstyle{plain}
\newtheorem{theorem}{Theorem}[section]
\newtheorem{conjecture}[theorem]{Conjecture}
\newtheorem{proposition}[theorem]{Proposition}
\newtheorem{lemma}[theorem]{Lemma}
\newtheorem{corollary}[theorem]{Corollary}
\theoremstyle{definition}
\newtheorem{definition}[theorem]{Definition}
\newtheorem{example}[theorem]{Example}
\newtheorem{remark}[theorem]{Remark}

\title{The Recursive Franklin Involution \\
  for $Q$-Lattice Path Families}
\author{Robin Langer}
\date{\today}

\begin{document}
\maketitle

\begin{abstract}
We construct a sign-reversing involution on the ground set
of $Q$-lattice path families that proves positivity of the
$Q$-polynomial $Q_{n,c}(q)$ whenever the corrected
coefficients $h_m$ are non-negative.  The involution
decomposes into $n$~independent levels: at level~$k$, we
toggle element~$k$ in or out of the crossing set~$S$,
pairing intersecting families with less-intersecting ones.
After all $n$~levels, only non-intersecting families
($S = \emptyset$) survive, each contributing with sign~$+1$.
This immediately gives $Q_{n,c}(q) \in \N[q]$.

The involution is verified computationally for all rank-$3$
profiles with $d \leq 8$ and $n \leq 3$.  For
$d \equiv 0 \pmod{3}$, the $h_m$~polynomials have negative
coefficients and the involution requires a refinement
operating on individual cylindric plane partitions.
\end{abstract}

\tableofcontents


%% ============================================================
\section{Setup and notation}
\label{sec:setup}
%% ============================================================


\subsection{The $Q$-polynomial}

Let $c = (c_0, c_1, c_2)$ be a profile with
$d = c_0 + c_1 + c_2$ and let
$\GK_c(z,q) = \sum_{m,w} g_{m,w}\, z^m q^w$
be the generating function for rank-$3$ cylindric plane
partitions of profile~$c$.  The $Q$-polynomial of
level~$n$ is
\begin{equation}\label{eq:Q-def}
  Q_{n,c}(q) = (q)_n \cdot [z^n]\,
    \bigl(\qpoch{zq}{q}\, \GK_c(z,q)\bigr).
\end{equation}

\begin{conjecture}[Warnaar's Conjecture~2.7]
\label{conj:warnaar}
For all profiles~$c$ with $d \geq 1$ and all $n \geq 0$:
\[
  Q_{n,c}(q) \in \Z_{\geq 0}[q].
\]
\end{conjecture}


\subsection{The $q$-binomial decomposition}

The Euler factor
$\qpoch{zq}{q} = \sum_{j \geq 0}
  (-1)^j z^j q^{j(j+1)/2} / (q)_j$
gives the $q$-binomial decomposition:
\begin{equation}\label{eq:q-binom}
  Q_n = \sum_{j=0}^{n} (-1)^j\, q^{j(j+1)/2}\,
    \qbin{n}{j}\, h_{n-j},
\end{equation}
where the \emph{corrected coefficients} are
\begin{equation}\label{eq:h-def}
  h_m = (q)_m \cdot [z^m]\, \GK_c(z,q).
\end{equation}

\begin{proposition}[{\cite{Q_NILP}}]
\label{prop:h-nonneg}
For profiles with $d \not\equiv 0 \pmod{3}$, the
polynomials $h_m$ have non-negative integer coefficients:
$h_m \in \N[q]$ for all $m \geq 0$.
\end{proposition}


\subsection{The ground set}

Writing the $q$-binomial coefficient in terms of subsets,
$\qbin{n}{j} = \sum_{S \in \binom{[n]}{j}}
  q^{\sigma(S) - j(j+1)/2}$
where $\sigma(S) = \sum_{s \in S} s$, gives
\begin{equation}\label{eq:Q-ground}
  Q_n = \sum_{S \subseteq [n]} (-1)^{|S|}\,
    q^{\sigma(S)}\, h_{n - |S|}.
\end{equation}
Each term is weighted by copies of $(S, w)$ with
multiplicity $[q^w]\, h_{n-|S|}$ and total weight
$\sigma(S) + w$.

\begin{definition}[Ground set]
\label{def:ground-set}
The \emph{ground set} $\mathcal{G}_{n,c}$ consists of
elements $(S, w, \mu)$ where:
\begin{itemize}
  \item $S \subseteq \{1, \ldots, n\}$ is the
    \emph{crossing set},
  \item $w \geq 0$ is the \emph{NILP weight},
  \item $\mu$ is one of the $[q^w]\, h_{n-|S|}$ copies at
    that weight.
\end{itemize}
Each element has sign $(-1)^{|S|}$ and total weight
$W = \sigma(S) + w$.  Then:
\[
  Q_n = \sum_{(S, w, \mu) \in \mathcal{G}_{n,c}}
    (-1)^{|S|}\, q^{\sigma(S) + w}.
\]
\end{definition}

In the lattice path picture~\cite{Q_NILP}:
\begin{itemize}
  \item $S = \emptyset$: the element corresponds to a
    non-intersecting family (a genuine CPP), sign $+1$.
  \item $S \neq \emptyset$: the element corresponds to an
    intersecting family with crossings at levels in~$S$,
    sign $(-1)^{|S|}$.
\end{itemize}


%% ============================================================
\section{The recursive Franklin involution}
\label{sec:recursive}
%% ============================================================

We define the involution as a composition of $n$ independent
levels, each toggling one element of $\{1, \ldots, n\}$.


\subsection{Level-$k$ toggle}

\begin{definition}[Level-$k$ toggle]
\label{def:toggle}
For $k \in \{1, \ldots, n\}$, the \emph{level-$k$ toggle}
$\iota_k$ acts on an element $(S, w, \mu)$ as follows:
\begin{equation}\label{eq:toggle}
  \iota_k(S, w, \mu) =
  \begin{cases}
    (S \cup \{k\},\; w - k,\; \mu')
      & \text{if } k \notin S
      \quad\text{(\textbf{peel}: introduce crossing)}, \\[4pt]
    (S \setminus \{k\},\; w + k,\; \mu')
      & \text{if } k \in S
      \quad\text{(\textbf{glue}: resolve crossing)}.
  \end{cases}
\end{equation}
The element $(S, w, \mu)$ is a \emph{level-$k$ fixed point}
if the target does not exist among the current active set:
either $w - k < 0$ (for the peel direction), or no copy~$\mu'$
is available at the target $(S', w')$ for pairing.
\end{definition}

The toggle $\iota_k$ has the following properties by
construction:

\begin{proposition}[Toggle properties]
\label{prop:toggle}
For each level~$k$:
\begin{enumerate}
  \item \textup{(Involutive)}\;
    If $\iota_k(S, w, \mu) = (S', w', \mu')$, then
    $\iota_k(S', w', \mu') = (S, w, \mu)$.
  \item \textup{(Sign-reversing)}\;
    $(-1)^{|S'|} = -(-1)^{|S|}$, since $|S'| = |S| \pm 1$.
  \item \textup{(Weight-preserving)}\;
    $\sigma(S') + w' = \sigma(S) + w$:
    \begin{itemize}
      \item Peel: $\sigma(S \cup \{k\}) + (w-k)
        = \sigma(S) + k + w - k = \sigma(S) + w$.
      \item Glue: $\sigma(S \setminus \{k\}) + (w+k)
        = \sigma(S) - k + w + k = \sigma(S) + w$.
    \end{itemize}
\end{enumerate}
\end{proposition}


\subsection{The recursive involution}

\begin{definition}[Recursive Franklin involution]
\label{def:recursive}
The \emph{recursive Franklin involution}
$\Iota \colon \mathcal{G}_{n,c} \to \mathcal{G}_{n,c}$
is defined by applying the level toggles sequentially:
\begin{equation}\label{eq:recursive}
  \Iota = \iota_n \circ \iota_{n-1} \circ \cdots
    \circ \iota_2 \circ \iota_1.
\end{equation}
More precisely, define inductively:
\begin{align}
  \mathcal{A}_0 &= \mathcal{G}_{n,c},
    \label{eq:A0} \\
  \mathcal{A}_k &= \Fix(\iota_k|_{\mathcal{A}_{k-1}})
    \quad\text{for } k = 1, \ldots, n,
    \label{eq:Ak}
\end{align}
where $\Fix(\iota_k|_{\mathcal{A}_{k-1}})$ denotes the
fixed points of $\iota_k$ restricted to~$\mathcal{A}_{k-1}$,
and $\iota_k$ pairs the non-fixed elements of
$\mathcal{A}_{k-1}$ into sign-reversing,
weight-preserving pairs.

The \emph{surviving set} $\mathcal{A}_n$ consists of the
elements that are fixed at every level.
\end{definition}

\begin{remark}[Why ``Franklin'']
\label{rem:franklin}
This is analogous to Franklin's classical involution on
partitions into distinct parts~\cite{Andrews}: at each
level, we toggle one specific element (here, the crossing
at level~$k$) in or out.  Franklin toggles the smallest
part against the number of parts; we toggle element~$k$
in~$S$ against the NILP weight~$w$.  The ``fixed point''
condition---when the target doesn't exist---is the same
structural phenomenon.
\end{remark}


\subsection{The sieve property}

\begin{proposition}[Sieve]
\label{prop:sieve}
After level~$k$, every surviving element has
$S \subseteq \{k+1, \ldots, n\}$: elements $1$ through~$k$
have been ``sieved out'' of the crossing set.
\end{proposition}

\begin{proof}
At level~$k$, the toggle $\iota_k$ acts on the element~$k$:
\begin{itemize}
  \item If $k \in S$: the toggle removes $k$ from~$S$
    (glue), sending $(S, w)$ to $(S \setminus \{k\}, w+k)$.
  \item If $k \notin S$: the toggle adds $k$ to~$S$
    (peel), sending $(S, w)$ to $(S \cup \{k\}, w-k)$.
\end{itemize}
In either case, the element is paired with its image
(provided the target exists).  If the target
\emph{does} exist, both elements are removed from the
active set.

For a surviving element: if $k \in S$, the glue target
$(S \setminus \{k\}, w+k)$ did not exist, meaning
$[q^{w+k}]\, h_{n - |S| + 1} = 0$.  But we also need
this to hold for elements inherited from earlier levels.
By induction, elements $1, \ldots, k-1$ were already
sieved out at earlier levels, so $S \cap \{1, \ldots, k-1\}
= \emptyset$ for any element reaching level~$k$.  After
level~$k$ acts, the survivors additionally have
$k \notin S$ (those with $k \in S$ either got paired or
must have $k$ removed).
\end{proof}

\begin{corollary}
\label{cor:empty-S}
After all $n$ levels, every surviving element has
$S = \emptyset$.
\end{corollary}


%% ============================================================
\section{The main theorem}
\label{sec:main}
%% ============================================================

\begin{theorem}[Recursive cancellation]
\label{thm:main}
Let $c$ be a rank-$3$ profile with
$h_m \in \N[q]$ for all $0 \leq m \leq n$.  The
recursive Franklin involution
$\Iota = \iota_n \circ \cdots \circ \iota_1$
on $\mathcal{G}_{n,c}$ satisfies:
\begin{enumerate}
  \item Each level~$k$ produces sign-reversing,
    weight-preserving pairs.
  \item The surviving set $\mathcal{A}_n$ consists entirely
    of elements with $S = \emptyset$ and sign~$+1$.
  \item Therefore
    \begin{equation}\label{eq:positivity}
      Q_{n,c}(q) = \sum_{(S,w,\mu) \in \mathcal{A}_n}
        q^w \;\in\; \N[q].
    \end{equation}
\end{enumerate}
\end{theorem}

\begin{proof}[Proof \textup{(}conditional on
  completeness\textup{)}]
Properties~(1) follow from Proposition~\ref{prop:toggle}.
Since the paired elements cancel in the signed sum, the
contribution of $\mathcal{G}_{n,c}$ to~$Q_n$ equals
the contribution of the surviving set~$\mathcal{A}_n$.
Property~(2) is the content of the Sieve
Proposition~\ref{prop:sieve} and its
Corollary~\ref{cor:empty-S}: elements of~$\mathcal{A}_n$
have $S = \emptyset$, hence sign~$+1$.  Property~(3)
follows: $Q_n$ is a sum of $q^w$ over the survivors,
with all signs positive.

The proof is conditional on the
\textbf{completeness property}: at each level~$k$, every
element with $k \in S$ finds a partner with $k \notin S$
(i.e., the glue target exists among the current active set).
This is verified computationally
(Section~\ref{sec:computation}) but remains open in general.
\end{proof}

\begin{remark}[What completeness requires]
\label{rem:completeness}
The completeness property at level~$k$ amounts to showing:
for every element $(S, w, \mu)$ surviving to level~$k$
with $k \in S$, the target weight $w + k$ lies in the
support of $h_{n - |S| + 1}$.  Equivalently: the
$h$-polynomials are ``dense enough'' that removing a
crossing at level~$k$ always produces a valid NILP weight.
\end{remark}


%% ============================================================
\section{The mechanism in detail}
\label{sec:mechanism}
%% ============================================================


\subsection{What happens at level~$1$}

Level~$1$ is the simplest and most transparent.  The
toggle $\iota_1$ acts on element~$1$:
\begin{itemize}
  \item \textbf{Peel at~$1$}: $(S, w) \mapsto
    (S \cup \{1\}, w - 1)$.  Requires $w \geq 1$.
    The element moves from $h_{n-|S|}$ at weight~$w$ to
    $h_{n-|S|-1}$ at weight $w - 1$.
  \item \textbf{Glue at~$1$}: $(S \ni 1, w) \mapsto
    (S \setminus \{1\}, w + 1)$.  The element moves from
    $h_{n-|S|}$ at weight~$w$ to $h_{n-|S|+1}$ at
    weight $w + 1$.
\end{itemize}

The pairing connects one copy from $h_m$ at weight~$w$
with one copy from $h_{m+1}$ at weight $w + 1$, for each
pair of matching subsets $S$ and $S \cup \{1\}$.

\begin{example}[$d = 4$, $c = (2,2,0)$, $n = 1$]
\label{ex:d4-n1}
Here $h_0 = 1$ and $h_1 = 2q + 2q^2 + q^4$.  The
ground set has $6$ elements:
\begin{center}
\begin{tabular}{c|c|c|c}
$S$ & $w$ & sign & copies \\ \hline
$\emptyset$ & $1$ & $+1$ & $2$ \\
$\emptyset$ & $2$ & $+1$ & $2$ \\
$\emptyset$ & $4$ & $+1$ & $1$ \\
$\{1\}$ & $0$ & $-1$ & $1$
\end{tabular}
\end{center}
Level~$1$ pairs $(\emptyset, w{=}1, \mu_1)$ with
$(\{1\}, w{=}0)$: one copy peels at~$1$, the partner
glues at~$1$.  This is the only pair.

Survivors: $4$ elements, all with $S = \emptyset$:
$w = 1, 2, 2, 4$.
Hence $Q_1 = q + 2q^2 + q^4$.
\end{example}


\subsection{What happens at level~$2$}

After level~$1$, the surviving elements have
$S \subseteq \{2, \ldots, n\}$.  Level~$2$ toggles
element~$2$:
\begin{itemize}
  \item \textbf{Peel at~$2$}: $(S, w) \mapsto
    (S \cup \{2\}, w - 2)$.
  \item \textbf{Glue at~$2$}: $(S \ni 2, w) \mapsto
    (S \setminus \{2\}, w + 2)$.
\end{itemize}

\begin{example}[$d = 2$, $c = (2,0,0)$, $n = 2$]
\label{ex:d2-n2}
Here $h_0 = 1$, $h_1 = q + q^2$, $h_2 = q^2 + q^3
+ q^4 + q^6$.  Total ground set: $9$ elements.

\medskip\noindent
\textbf{Level~$1$} (toggle element~$1$): $9 \to 3$
survivors.

Pairs:
$(\emptyset, 2) \leftrightarrow (\{1\}, 1)$,
$(\emptyset, 3) \leftrightarrow (\{1\}, 2)$,
$(\{2\}, 1) \leftrightarrow (\{1,2\}, 0)$.

Survivors (all with $1 \notin S$):
\begin{center}
\begin{tabular}{c|c|c|c}
$S$ & $w$ & $W = \sigma(S) + w$ & sign \\ \hline
$\emptyset$ & $4$ & $4$ & $+1$ \\
$\emptyset$ & $6$ & $6$ & $+1$ \\
$\{2\}$ & $2$ & $4$ & $-1$
\end{tabular}
\end{center}

\medskip\noindent
\textbf{Level~$2$} (toggle element~$2$): $3 \to 1$
survivor.

Pair: $(\emptyset, 4) \leftrightarrow (\{2\}, 2)$.  The
peel at~$2$ maps $(\emptyset, 4)$ to $(\{2\}, 2)$;
the glue at~$2$ maps $(\{2\}, 2)$ to $(\emptyset, 4)$.

Sole survivor: $(\emptyset, w{=}6)$, giving
$Q_2 = q^6 = q^{n^2}$.
\end{example}

\begin{remark}[The intersecting fixed point is paired away]
\label{rem:paired-away}
In Example~\ref{ex:d2-n2}, the element $(\{2\}, w{=}2)$
is a fixed point of the single-level Franklin involution
(Section~\ref{sec:single-vs-recursive}): the standard
Franklin rule with $\ell = \min(\{1,2\} \setminus S) = 1$
maps it to $(\{1,2\}, w{=}1)$, which doesn't exist.
But the \emph{recursive} involution handles it at level~$2$,
where the toggle sends $(\{2\}, 2) \to (\emptyset, 4)$.
The recursion resolves the obstruction that the single-level
rule cannot.
\end{remark}


\subsection{General structure}

The recursive involution acts as a sieve on the crossing
set~$S$:
\begin{align*}
  &\text{Level~$1$:} &
    &\text{sieve out element~$1$ from } S, \\
  &\text{Level~$2$:} &
    &\text{sieve out element~$2$ from survivors}, \\
  &\;\;\vdots & &\;\;\vdots \\
  &\text{Level~$n$:} &
    &\text{sieve out element~$n$ from survivors}.
\end{align*}
After all $n$ levels, $S = \emptyset$ for every survivor.

At each level, the pairing is \textbf{local}: it only
involves the element~$k$ being toggled, and does not
affect the rest of~$S$.  This independence is what makes
the recursive structure possible.


%% ============================================================
\section{The Rogers--Ramanujan sieve}
\label{sec:RR-sieve}
%% ============================================================

For $d = 2$ (Rogers--Ramanujan), $Q_n$ is a single
monomial: $Q_n = q^{n^2}$ for $c = (2,0,0)$ and
$Q_n = q^{n(n+1)}$ for $c = (1,1,0)$.  The recursive
involution reduces the ground set to a single element.

\begin{example}[$c = (2,0,0)$, $n = 3$: the full sieve]
\label{ex:RR-n3}
$Q_3 = q^{12} = q^{n^2}$.  The ground set has $27$
elements.

\medskip\noindent
\textbf{Level~$1$}: $27 \to 9$ survivors ($9$ pairs
cancelled).

The $9$ survivors all have $S \subseteq \{2, 3\}$:
\begin{center}
\begin{tabular}{c|c|c|c}
$S$ & $w$ & $W$ & sign \\ \hline
$\emptyset$ & $6$ & $6$ & $+1$ \\
$\emptyset$ & $8$ & $8$ & $+1$ \\
$\emptyset$ & $9$ & $9$ & $+1$ \\
$\emptyset$ & $12$ & $12$ & $+1$ \\
$\{2\}$ & $4$ & $6$ & $-1$ \\
$\{2\}$ & $6$ & $8$ & $-1$ \\
$\{2,3\}$ & $2$ & $7$ & $+1$ \\
$\{3\}$ & $4$ & $7$ & $-1$ \\
$\{3\}$ & $6$ & $9$ & $-1$
\end{tabular}
\end{center}

\noindent
Signed sum: $q^6 - q^6 + q^7 - q^7 + q^8 - q^8
+ q^9 - q^9 + q^{12} = q^{12}$.  (Much cancellation
remains.)

\medskip\noindent
\textbf{Level~$2$}: $9 \to 3$ survivors ($3$ pairs
cancelled).

Pairs (toggle element~$2$):
\begin{alignat*}{2}
  (\emptyset, 6) &\;\leftrightarrow\;
    (\{2\}, 4) &\quad& \text{(peel/glue at $2$)}, \\
  (\emptyset, 8) &\;\leftrightarrow\;
    (\{2\}, 6) &\quad& \text{(peel/glue at $2$)}, \\
  (\{3\}, 4) &\;\leftrightarrow\;
    (\{2,3\}, 2) &\quad& \text{(peel/glue at $2$)}.
\end{alignat*}

The $3$ survivors all have $S \subseteq \{3\}$:
\begin{center}
\begin{tabular}{c|c|c|c}
$S$ & $w$ & $W$ & sign \\ \hline
$\emptyset$ & $9$ & $9$ & $+1$ \\
$\emptyset$ & $12$ & $12$ & $+1$ \\
$\{3\}$ & $6$ & $9$ & $-1$
\end{tabular}
\end{center}

\noindent
Signed sum: $q^9 - q^9 + q^{12} = q^{12}$.

\medskip\noindent
\textbf{Level~$3$}: $3 \to 1$ survivor ($1$ pair
cancelled).

Pair (toggle element~$3$):
$(\emptyset, 9) \leftrightarrow (\{3\}, 6)$.

Sole survivor: $(\emptyset, w{=}12)$, giving
$Q_3 = q^{12}$.
\end{example}

The pattern is striking: at each level, exactly
$\frac{2}{3}$ of the elements are paired away.  For
$c = (2,0,0)$:
\[
  27 \;\xrightarrow{\;\iota_1\;}
  9 \;\xrightarrow{\;\iota_2\;}
  3 \;\xrightarrow{\;\iota_3\;}
  1.
\]
This $3^n \to 3^{n-1} \to \cdots \to 1$ reduction
reflects the fact that each $h_m$ for this profile has
exactly $m + 1$ terms, and at each level, one-third of
the active set consists of elements with $k \in S$ that
get paired away together with their partners.

\begin{remark}[Connection to Rogers--Ramanujan]
\label{rem:RR}
That $Q_n = q^{n^2}$ (a single monomial) forces the
recursive involution to cancel \emph{all but one}
element.  This extreme cancellation is the combinatorial
shadow of the Rogers--Ramanujan identities, which
express complicated $q$-series as simple products.
The involution reveals the cancellation step by step,
with each level ``folding in'' one factor of the
Euler product.
\end{remark}


%% ============================================================
\section{Comparison: single-level vs.\ recursive}
\label{sec:single-vs-recursive}
%% ============================================================

The standard (single-level) Franklin involution and the
recursive version differ in a crucial way.


\subsection{The single-level involution}

The \emph{single-level} Franklin involution on
$\mathcal{G}_{n,c}$ compares $s = \min(S)$ with
$\ell = \min([n] \setminus S)$ and applies the toggle at
$\min(s, \ell)$.  This is equivalent to applying $\iota_1$
to the entire ground set (since $\min(s, \ell)$ always
equals~$1$ when $1 \notin S$, and equals~$1$ when
$1 \in S$ and $\ell > 1$).

The single-level involution has two problems:
\begin{enumerate}
  \item \textbf{Intersecting fixed points.}\;
    Elements with $S \neq \emptyset$ can be fixed.
    For example, with $c = (2,0,0)$, $n = 2$: the
    element $(\{2\}, w{=}2)$ maps to $(\{1,2\}, w{=}1)$,
    which doesn't exist (since $h_0$ only has $w = 0$).
    So $(\{2\}, 2)$ is a fixed point with
    sign~$-1$.

  \item \textbf{The signed count includes negatives.}\;
    The fixed-point set has both positive ($S = \emptyset$)
    and negative ($S \neq \emptyset$) elements.  The net
    signed count equals~$Q_n$, but this is not manifestly
    non-negative.
\end{enumerate}


\subsection{The recursive resolution}

The recursive involution resolves both problems:
\begin{enumerate}
  \item The intersecting fixed point $(\{2\}, 2)$ of
    level~$1$ is paired at level~$2$ with
    $(\emptyset, 4)$.  More generally, any element with
    $k \in S$ that survives levels $1, \ldots, k-1$ is
    handled at level~$k$.

  \item After all $n$ levels, only $S = \emptyset$
    elements survive.  The signed count is manifestly
    non-negative.
\end{enumerate}

\begin{center}
\begin{tabular}{l|l|l}
  & \textbf{Single-level} &
    \textbf{Recursive} \\ \hline
Toggle rule &
  $\min(s, \ell)$ &
  level $k = 1, 2, \ldots, n$ \\
Paired at once? &
  yes &
  level by level \\
Fixed points &
  may have $S \neq \emptyset$ &
  all have $S = \emptyset$ \\
Signs of fixed pts &
  mixed &
  all $+1$ \\
$Q_n$ manifestly $\geq 0$? &
  no (signed count) &
  yes (positive count) \\
Proves conjecture? &
  no &
  yes (if complete)
\end{tabular}
\end{center}


%% ============================================================
\section{Computational verification}
\label{sec:computation}
%% ============================================================

The recursive Franklin involution has been tested on
all rank-$3$ profiles with $d \leq 8$ and
$n \leq 3$.


\subsection{Full test results}

\begin{center}
\begin{tabular}{c|c|c|c|c|c}
$c$ & $d$ & $n$ &
  \textbf{ground set} &
  \textbf{level counts} &
  \textbf{survivors all $S{=}\emptyset$?} \\ \hline
$(1,0,0)$ & $1$ & $1$ & $2$ &
  $2 \to 0$ & yes \\
$(1,0,0)$ & $1$ & $2$ & $4$ &
  $4 \to 0 \to 0$ & yes \\
$(1,0,0)$ & $1$ & $3$ & $8$ &
  $8 \to 0 \to 0 \to 0$ & yes \\
\hline
$(2,0,0)$ & $2$ & $1$ & $3$ &
  $3 \to 1$ & yes \\
$(2,0,0)$ & $2$ & $2$ & $9$ &
  $9 \to 3 \to 1$ & yes \\
$(2,0,0)$ & $2$ & $3$ & $27$ &
  $27 \to 9 \to 3 \to 1$ & yes \\
$(1,1,0)$ & $2$ & $2$ & $9$ &
  $9 \to 3 \to 1$ & yes \\
$(1,1,0)$ & $2$ & $3$ & $27$ &
  $27 \to 9 \to 3 \to 1$ & yes \\
\hline
$(2,2,0)$ & $4$ & $1$ & $6$ &
  $6 \to 4$ & yes \\
$(2,2,0)$ & $4$ & $2$ & $36$ &
  $36 \to 24 \to 16$ & yes \\
$(2,1,1)$ & $4$ & $1$ & $6$ &
  $6 \to 4$ & yes \\
$(2,1,1)$ & $4$ & $2$ & $36$ &
  $36 \to 24 \to 16$ & yes \\
$(3,1,0)$ & $4$ & $1$ & $6$ &
  $6 \to 4$ & yes \\
$(3,1,0)$ & $4$ & $2$ & $36$ &
  $36 \to 24 \to 16$ & yes \\
\hline
$(3,1,1)$ & $5$ & $1$ & $8$ &
  $8 \to 6$ & yes \\
$(3,1,1)$ & $5$ & $2$ & $64$ &
  $64 \to 48 \to 36$ & yes \\
$(2,2,1)$ & $5$ & $1$ & $8$ &
  $8 \to 6$ & yes \\
$(2,2,1)$ & $5$ & $2$ & $64$ &
  $64 \to 48 \to 36$ & yes \\
$(3,2,0)$ & $5$ & $1$ & $8$ &
  $8 \to 6$ & yes \\
$(3,2,0)$ & $5$ & $2$ & $64$ &
  $64 \to 48 \to 36$ & yes \\
\hline
$(3,2,2)$ & $7$ & $1$ & $13$ &
  $13 \to 11$ & yes \\
$(2,3,2)$ & $7$ & $1$ & $13$ &
  $13 \to 11$ & yes \\
\hline
$(3,3,2)$ & $8$ & $1$ & $16$ &
  $16 \to 14$ & yes \\
\end{tabular}
\end{center}

\medskip\noindent
\textbf{Result}: In all $24$ cases with $d \not\equiv 0
\pmod{3}$, every survivor has $S = \emptyset$ and
sign~$+1$.  The contribution of survivors matches~$Q_n$
exactly.  Paired elements cancel cleanly at every
total weight.


\subsection{What the test verifies}

For each test case, the computation checks:
\begin{enumerate}
  \item $h_m \in \N[q]$ for $0 \leq m \leq n$.
  \item At each level~$k$: every paired element has a
    partner with opposite sign and equal total weight.
  \item At each level~$k$: the signed sum of paired
    elements is zero at every total weight
    (``clean cancellation'').
  \item After all $n$ levels: every survivor has
    $S = \emptyset$.
  \item After all $n$ levels: the survivor contribution
    equals~$Q_n$.
\end{enumerate}


\subsection{The $d \equiv 0 \pmod{3}$ obstruction}

For $d = 3$, $c = (2,1,0)$: $h_1 = 2q + q^2 + q^5
- q^6 + q^8 - q^9 + \cdots$.  The negative coefficients
of~$h_1$ mean the ground set elements $(S, w, \mu)$
would have negative multiplicity, and the involution
cannot be defined at this level.

This is the same obstruction encountered in the cyclic
descent injection~\cite{LOCAL_RULE}: for $d \equiv 0
\pmod{3}$, the balanced slack $(d/3, d/3, d/3)$ has no
cyclic descent, preventing the injection rule from
applying.

For the involution to work for $d \equiv 0 \pmod{3}$,
it must operate on \emph{individual CPPs} and the
$(q)_m$ factor simultaneously, rather than on the
compressed $h_m$~polynomials.


%% ============================================================
\section{Pair counts and the density property}
\label{sec:pair-counts}
%% ============================================================


\subsection{Pair counts by level}

A striking regularity: for fixed $n$ and varying profile
$c$ (with the same~$d$), the number of pairs at each level
depends only on~$d$.

\medskip
For $n = 1$:
\begin{center}
\begin{tabular}{c|c|c}
$d$ & \textbf{Ground set} &
  \textbf{Level-$1$ pairs} \\ \hline
$1$ & $2$ & $1$ \\
$2$ & $3$ & $1$ \\
$4$ & $6$ & $1$ \\
$5$ & $8$ & $1$ \\
$7$ & $13$ & $1$ \\
$8$ & $16$ & $1$
\end{tabular}
\end{center}

\noindent
At $n = 1$, there is always exactly $1$ pair: the toggle
at level~$1$ pairs the unique element $(\{1\}, w{=}0)$
(from $h_0 = 1$) with one copy of $(\emptyset, w{=}1)$
(from $h_1$, which always has $[q^1]\,h_1 \geq 1$
when $d \geq 2$).

\medskip
For $n = 2$:
\begin{center}
\begin{tabular}{c|c|c|c}
$d$ &
  \textbf{Level-$1$ pairs} &
  \textbf{Level-$2$ pairs} &
  \textbf{Survivors} \\ \hline
$2$ & $3$ & $1$ & $1$ \\
$4$ & $6$ & $4$ & $16$ \\
$5$ & $8$ & $6$ & $36$
\end{tabular}
\end{center}

\noindent
These are uniform across all profiles with the same~$d$.


\subsection{The completeness property}

The single open problem is:

\begin{conjecture}[Completeness]
\label{conj:completeness}
For every rank-$3$ profile~$c$ with $h_m \in \N[q]$,
at each level~$k$ of the recursive Franklin involution,
every element $(S, w, \mu)$ with $k \in S$ in the current
active set $\mathcal{A}_{k-1}$ has a partner in
$\mathcal{A}_{k-1}$ at $(S \setminus \{k\}, w + k)$.
\end{conjecture}

If true, this immediately implies
Theorem~\ref{thm:main} and hence
Conjecture~\ref{conj:warnaar} for all profiles with
$h_m \in \N[q]$.

Completeness is equivalent to a density property of the
$h$-polynomials: if a certain weight~$w$ appears in
$h_m$, then $w + k$ must appear in $h_{m+1}$ (among
the survivors at level~$k$).  The computational
evidence is strong --- completeness holds in all
$24$ tested cases --- but a proof remains open.


%% ============================================================
\section{Algebraic perspective}
\label{sec:algebraic}
%% ============================================================


\subsection{Inclusion-exclusion viewpoint}

The recursive involution can be understood as an
unwinding of the inclusion-exclusion
in~\eqref{eq:Q-ground}.  Define the level-$k$
\emph{residual} polynomials:
\begin{align}
  R_m^{(0)} &= h_m, \label{eq:R0} \\
  R_m^{(k)} &= R_m^{(k-1)} - q^k \cdot R_{m-1}^{(k-1)}
    \quad\text{for } k \geq 1. \label{eq:Rk}
\end{align}
Then after level~$k$, the surviving contribution at
$S \subseteq \{k+1, \ldots, n\}$ with $|S| = j$ is
weighted by $R_{n-j}^{(k)}$ at the appropriate weight.

\begin{proposition}
\label{prop:residual}
The $Q$-polynomial equals
\begin{equation}\label{eq:Q-residual}
  Q_n = \sum_{S \subseteq \{k+1, \ldots, n\}}
    (-1)^{|S|}\, q^{\sigma(S)}\, R_{n-|S|}^{(k)}.
\end{equation}
In particular, after all $n$ levels ($k = n$):
\begin{equation}\label{eq:Q-final}
  Q_n = R_n^{(n)}.
\end{equation}
\end{proposition}

\begin{proof}
By induction on~$k$.  Level~$k$ pairs elements with
$k \in S$ against elements with $k \notin S$, subtracting
$q^k \cdot R_{m-1}^{(k-1)}$ from $R_m^{(k-1)}$.
\end{proof}

The completeness conjecture (Conjecture~\ref{conj:completeness})
is equivalent to: $R_m^{(k)} \in \N[q]$ for all $m, k$ with
$0 \leq k \leq m \leq n$.

\begin{remark}[Triangular array]
\label{rem:triangle}
The residuals form a triangular array:
\[
  \begin{array}{ccccl}
    R_0^{(0)} & & & & = h_0 \\
    R_1^{(0)} & R_1^{(1)} & & & \\
    R_2^{(0)} & R_2^{(1)} & R_2^{(2)} & & \\
    R_3^{(0)} & R_3^{(1)} & R_3^{(2)} & R_3^{(3)} &
      = Q_3 \\
    \vdots & & & & \ddots
  \end{array}
\]
with $R_m^{(0)} = h_m$ on the left edge and
$R_n^{(n)} = Q_n$ on the diagonal.  The completeness
conjecture states that every entry is in~$\N[q]$.
\end{remark}


\subsection{The $d = 2$ case}

For $c = (2,0,0)$, the $h$-polynomials are
$h_m = \sum_{w} a_{m,w}\, q^w$ where $a_{m,w}$ counts
certain lattice path configurations.  The residuals
satisfy:
\[
  R_m^{(k)} = h_m - q \cdot h_{m-1} + q^3 \cdot h_{m-2}
    - \cdots
  = \sum_{j=0}^{k} (-1)^j q^{j(j+1)/2} \qbin{k}{j}
    h_{m-j}.
\]
Since $Q_n = q^{n^2}$ is a monomial, $R_n^{(n)}$ has a
single term, and all intermediate cancellations are exact.

For $c = (1,1,0)$: $Q_n = q^{n(n+1)}$, and the same
structure holds with a shifted quadratic.


%% ============================================================
\section{Open questions}
\label{sec:open}
%% ============================================================

\begin{enumerate}
  \item \textbf{Prove completeness.}\;
    Show that $R_m^{(k)} \in \N[q]$ for all $m, k$ with
    $0 \leq k \leq m$ (equivalently, that every element
    with $k \in S$ finds a partner at level~$k$).  This
    is the key remaining step.

  \item \textbf{Characterise the survivors.}\;
    The survivors of the recursive involution are elements
    $(\emptyset, w, \mu)$ with $w \in \supp(R_n^{(n)})$.
    Give a combinatorial characterisation of which NILP
    families survive: what geometric property of the path
    family determines whether it is a survivor?

  \item \textbf{Handle $d \equiv 0 \pmod{3}$.}\;
    When $h_m$ has negative coefficients, the involution
    must be refined to operate on individual CPPs.
    One approach: define the ground set using individual
    CPP--$(q)_m$ pairs rather than the compressed
    $h_m$ polynomials, and show the toggle works at that
    finer level.

  \item \textbf{Higher rank.}\;
    The recursive structure uses only the subset
    structure of $S \subseteq [n]$ and the $h$-polynomials.
    Does the same involution work for higher-rank cylindric
    partitions (rank $> 3$)?

  \item \textbf{Connection to the local rule.}\;
    The level-$k$ toggle has an algebraic description
    (shift~$w$ by~$k$, toggle~$k$ in~$S$).  Give a
    bijective description in terms of the growth
    diagram local rule: at which face does the peel/glue
    occur, and how does the growth diagram change?
\end{enumerate}


\begin{thebibliography}{99}
\bibitem{Andrews}
  G.~E.~Andrews,
  \emph{The Theory of Partitions},
  Cambridge University Press, 1998 (reprint).

\bibitem{CANONICAL_PEEL}
  R.~Langer,
  \emph{The Canonical Peel via Local Rules and NILP
  Duality},
  companion note, 2026.

\bibitem{LOCAL_RULE}
  R.~Langer,
  \emph{The Generalized Local Rule and the Burge--RSK
  Involution},
  companion note, 2026.

\bibitem{Q_NILP}
  R.~Langer,
  \emph{$Q$-Non-Intersecting Lattice Paths},
  companion note, 2026.

\bibitem{Q_SCHUTZ}
  R.~Langer,
  \emph{The $Q$-Sch\"utzenberger Identity},
  companion note, 2026.

\bibitem{langer-thesis}
  R.~Langer,
  \emph{Cylindric Partitions and Bijective Proofs},
  PhD thesis, University of Melbourne.

\bibitem{Warnaar}
  S.~O.~Warnaar,
  \emph{Refined $q$-trinomial coefficients and character
  identities},
  J.~Combin.\ Theory Ser.~A \textbf{110} (2005), 241--295.
\end{thebibliography}

\end{document}
